\section{Implementación del software}

El análisis transversal de los 20 casos de estudio revela que la industria ha superado la etapa de experimentación con EDA para consolidar un \textbf{stack tecnológico de referencia}. Identificamos que el éxito de estas implementaciones no radica en una sola herramienta, sino en la orquestación precisa entre tres capas lógicas: la ingesta en el borde (Edge), el transporte asíncrono y el procesamiento distribuido.

A continuación, realizamos una deconstrucción profunda de las decisiones de ingeniería críticas que permitieron a sistemas como el de tráfico de Sri Lanka \cite{perera2024traffic} o el de monitoreo de Tecgraf \cite{laigner2020monolithic} escalar donde sus predecesores fallaron.

\subsection{El Stack de Visión y Detección (Edge)}
En la capa de captura, observamos una homogeneización en torno a soluciones de código abierto optimizadas para hardware limitado. La implementación descrita por Perera et al. \cite{perera2024traffic} es paradigmática en cuanto al manejo de concurrencia: utilizan un servidor Edge ejecutando scripts de \textbf{Python} con arquitectura \textit{Multi-threading}. Mientras un hilo dedicado captura el flujo de video, otro ejecuta la inferencia utilizando \textbf{YOLOv5} para la detección de objetos y \textbf{DeepSORT} para el rastreo (tracking) de trayectorias, asegurando que el bloqueo en la inferencia no congele la captura de frames.

Esta combinación es crucial. Como señalan Kul y Tashiev \cite{kul2021microservices}, delegar la extracción de características (tipo, color, velocidad) al borde utilizando redes como YOLOv4-tiny permite que lo que viaje por la red no sea video pesado, sino metadatos ligeros (JSON). En casos extremos de eficiencia, Coretti, Varile y Bertaina \cite{coretti2025neuromorphic} reemplazan el procesamiento de cuadros convencional por algoritmos \textbf{Stack-CNN}, que procesan únicamente los eventos de cambio de polaridad en sensores neuromórficos, eliminando la redundancia visual desde la raíz.

\subsubsection{Estrategias de Preprocesamiento Local}
Un hallazgo técnico relevante en la implementación de Perera es el uso de técnicas de visión clásica antes de invocar a la IA. Para la detección nocturna, implementan un pipeline de preprocesamiento que incluye conversión a escala de grises, umbralización adaptativa y operaciones morfológicas para eliminar el ruido de los faros de los coches antes de intentar leer la matrícula. Esto demuestra que la implementación robusta no confía ciegamente en la red neuronal, sino que "limpia" la entrada para maximizar la probabilidad de éxito.

\subsection{El Sistema Nervioso: Patrones de Mensajería}
Una vez generado el evento, la estrategia de transporte define la resiliencia del sistema. La literatura muestra dos enfoques dominantes dependiendo de la criticidad y el volumen de datos:

\begin{enumerate}
    \item \textbf{Bus de Eventos Estricto (Kafka):} En entornos de alta densidad de datos, como el análisis de video distribuido, Kul y Tashiev \cite{kul2021microservices} implementan \textbf{Apache Kafka} sobre contenedores \textbf{Docker}. Su innovación radica en el patrón de "topics jerárquicos" y "Mergers": un microservicio publica el color, otro la velocidad, y un tercero ("Merger") se suscribe a ambos para fusionar la información en un evento complejo (Type-Color-Speed). Esto garantiza un desacoplamiento total y permite que cada equipo de desarrollo evolucione sus microservicios sin romper a los demás.
    \item \textbf{Triggers Serverless (Firebase):} Para escenarios más orientados a la gestión administrativa, como el sistema de multas de Perera et al. \cite{perera2024traffic}, se opta por arquitecturas basadas en disparadores. La inserción de una infracción en \textbf{Firebase Realtime Database} dispara automáticamente una \textbf{Cloud Function} (HTTP Trigger) que notifica al servidor central. Esto elimina la necesidad de mantener un servidor de mensajería dedicado para eventos de baja frecuencia, optimizando costos operativos.
\end{enumerate}

\subsubsection{Protocolos de Ingesta}
Más allá de Kafka, Park et al. \cite{park2021cloud} detallan la importancia de la capa de traducción de protocolos. En su implementación de REDA, los sensores de bajo consumo no hablan Kafka nativamente (que usa TCP pesado); en su lugar, comunican vía \textbf{MQTT}, un protocolo ligero ideal para redes inestables. Un componente puente ("MQTT-Kafka Bridge") se encarga de ingerir estos mensajes y producirlos en el bus central, actuando como un amortiguador de picos de carga.

\subsection{Persistencia, Orquestación y Consistencia}
La transición a microservicios exige repensar el almacenamiento. Laigner et al. \cite{laigner2020monolithic} documentan cómo el modelo relacional tradicional se convierte en un cuello de botella bajo carga masiva. Su solución para el sistema de monitoreo de flotas fue migrar la capa de persistencia a \textbf{Cassandra}, una base de datos NoSQL optimizada para escrituras rápidas, esencial para guardar la telemetría de miles de sensores sin bloqueos de tabla.

Sin embargo, la coreografía de eventos pura puede volverse inmanejable y difícil de rastrear. Adekunle \cite{adekunle2024financial} advierte que, en sistemas financieros o críticos, dejar que los servicios reaccionen libremente es riesgoso ("infierno de la coreografía"). Por ello, la implementación robusta suele incluir un motor de procesos (como \textbf{Camunda BPM}) que orquesta los flujos de Kafka. Este motor mantiene el estado de la transacción y asegura que si un evento de "infracción detectada" no recibe un evento de "multa pagada" en un tiempo determinado, se dispare una alerta de escalado o compensación.

\subsection{Fragmento de código: Productor de Eventos IoT}
El siguiente código conceptualiza la lógica de un nodo Edge basada en el patrón de filtrado descrito por Salamon et al. \cite{salamon2018gateway}. A diferencia de un sistema tonto que transmite todo lo que ve, este script implementa lógica de negocio local para decidir qué merece ser enviado al bus global, optimizando el ancho de banda.

\begin{lstlisting}[language=Python]
from kafka import KafkaProducer
import json
import time

# Configuracion del productor (Edge Node)
# Se conecta al cluster Kafka gestionado por Zookeeper
producer = KafkaProducer(
    bootstrap_servers=['kafka-broker:9092'],
    value_serializer=lambda v: json.dumps(v).encode('utf-8')
)

def procesar_sensor_trafico(sensor_id, datos_vision):
    """
    Simula la logica de 'Smart Filtering' en el borde.
    Solo emite eventos si la confianza de la IA (YOLO) es alta.
    """
    nivel_confianza = datos_vision.get('confidence', 0)
    tipo_infraccion = datos_vision.get('type')

    # Filtrado en origen: Reduccion de ancho de banda y ruido
    # como sugieren Salamon et al.
    if nivel_confianza > 0.85:
        evento = {
            "sensor_id": sensor_id,
            "timestamp": time.time(),
            "evento": "INFRACCION_DETECTADA",
            "detalles": {
                "tipo": tipo_infraccion,
                "confianza": nivel_confianza,
                "coordenadas": datos_vision.get('bbox')
            }
        }
        
        # Publicacion asincrona al topico particionado
        # La clave de particion podria ser el sensor_id
        producer.send('alertas-trafico', key=sensor_id, value=evento)
        print(f"Evento critico enviado: {tipo_infraccion}")
    else:
        # Los falsos positivos se descartan localmente
        # Esto reduce la carga en el consumidor central
        pass
\end{lstlisting}